\section{Co-PEER and BugBank: Reusable Mathematical Objections}
\label{sec:resource}
\label{sec:resource-audit}
Co-PEER is our workflow for constructing and checking the proofs of error defined in Section~\ref{sec:evidence}, and BugBank is the collection that stores each resulting objection dossier and its revision history for inspection and reuse.

\subsection{From a selected assertion to an objection dossier}
\label{sec:pipeline}
Given a problem statement and a numbered proof, Co-PEER constructs an objection dossier in six stages (prompts, search bounds, and procedures are given in Appendix~\ref{sec:method}; see also \citealp{mcconnell2011certifying,moura2021lean,mathlib2020}):
\begin{itemize}[leftmargin=*,itemsep=2pt,topsep=3pt]
\item \textbf{1. Selecting a candidate assertion.} Two calls to a fixed language model read the problem and the numbered proof and each select an assertion that may admit a short counterexample, quoting the source verbatim and stating its assumptions and conclusion. A deterministic validator checks the exact quotation and output format, and a second model reviews each candidate.
\item \textbf{2. Searching for a counterexample.} A search agent writes and runs exact-arithmetic programs to test the assertion. It records one of four outcomes: a concrete witness, no witness found after exhaustive search of a specified finite set, an inconclusive result, or a claim outside the permitted search setting.
\item \textbf{3. Checking the witness independently.} A separately written checker program reconstructs the condition from the quoted claim and tests the witness independently of the search script.
\item \textbf{4. Formalizing the certificate in Lean.} Human contributors and coding agents construct the formal certificate of Equation~\eqref{eq:formal-witness}; Lean checks the proof, and an axiom audit records the assumptions and dependencies used.
\item \textbf{5. Reviewing the mathematics and source correspondence.} Three fresh model reviews audit the dossier. One searches for a counterexample from the problem and quoted claim alone; one tests the recorded witness with its own program; and one compares the Lean certificate with the quoted claim and witness to assess Equation~\eqref{eq:transport}. A further model examines flagged records and recommends keeping, repairing, excluding, or referring them for human review.
\item \textbf{6. Recording the dossier.} The workflow packages the stage inputs, outputs, certificates, and any recorded review decisions into a persistent objection dossier.
\end{itemize}

\noindent\begin{minipage}{\linewidth}
Each \emph{objection dossier} preserves four groups of evidence for inspection:

{\small\noindent
\begin{tabularx}{\linewidth}{@{}p{.27\linewidth}X@{}}
\toprule
Dossier component & What a reader can inspect\\\midrule
Source and claim & The saved source version, quoted passage, selected assertion, and the assumptions under which that assertion is claimed to hold.\\
Counterexample and proof & The concrete witness, formal definitions, and exact named Lean theorem proving Equation~\eqref{eq:formal-witness}.\\
Checking record & File identifiers, dependencies, build logs, and the axioms and assumptions used by the formal proof.\\
Review and consequence & Assessments of source correspondence~\eqref{eq:transport} and of what the objection shows about the surrounding argument (local or global), with the reviewer (human or model), reviewed version, verdict, and any unresolved concerns recorded for each assessment.\\
\bottomrule
\end{tabularx}\par}
\end{minipage}
The dossier records what the objection shows about the surrounding argument separately from the checks of the counterexample and its correspondence to the selected source claim.

\paragraph{BugBank.}
BugBank collects these dossiers. A dossier in BugBank shows what the workflow produced and which reviews, if any, it received; human review is not a condition of membership in the broader collection. Different subsets were reviewed by humans for different studies. Section~\ref{sec:evaluation} describes those reviews and defines the 147-dossier \emph{human-approved core} used for the benchmark.

\paragraph{Why preserve the history?}
A proof checked by Lean does not by itself establish that an objection applies to the source. In the Shellsort example of Section~\ref{sec:evidence}, the initial witness $(3,4,2,5,1)$ had a proof checked by Lean, but source review found that it violates the source's endpoint assumption ($3>1$). The replacement $(3,5,2,6,1,7,4)$ satisfies that assumption ($3<4$), refutes the same step, and also has a Lean certificate. By keeping both versions and the reason for the replacement, the dossier lets a reader check which objection applies under the source's assumptions before reusing it. Section~\ref{app:source-admissibility} compares the two witnesses; Appendix~\ref{sec:method} and Section~\ref{app:verified-census} document construction and replay.
